\documentclass[10pt,UTF8]{article}

\usepackage{geometry}

\usepackage{ctex}

\geometry{a4paper,scale=0.9}
\ctexset{today=old}

\usepackage[bookmarksnumbered=true]{hyperref}  % 生成大纲


\usepackage{times}  % 设置正文字体为 Adobe Times Roman

% \usepackage{makecell}
% \usepackage{multirow} % 表格多行合并

% \usepackage{graphicx} %插入图片的宏包
% \usepackage{subfigure}
% \usepackage{float} %设置图片浮动位置的宏包

\usepackage{amsmath}
\usepackage{amssymb}
\usepackage{mathtools}

\usepackage{physics}

% \usepackage{siunitx}
% % 关闭警告
% \ExplSyntaxOn
% \msg_redirect_name:nnn { siunitx } { physics-pkg } { none }
% \ExplSyntaxOff

%%%%%%%%%%%%%%%%%%%%%%%% tcolorbox %%%%%%%%%%%%%%%%%%%%%%%%%%
\usepackage[most]{tcolorbox}

\newenvironment{definition box}
{\begin{tcolorbox}[colback=blue!5!white,colframe=blue!75!black,parbox=false,before upper=\par,before lower=\par]}
{\end{tcolorbox}}

\newenvironment{theorem box}
{\begin{tcolorbox}[colback=yellow!5!white,colframe=yellow!75!black,parbox=false,before upper=\par,before lower=\par]}
{\end{tcolorbox}}
%%%%%%%%%%%%%%%%%%%%%%%%%%%%%%%%%%%%%%%%%%%%%%%%%%%%%%%%%%%%
            
%%%%%%%%%%%%%%%%%%%%% amsthm %%%%%%%%%%%%%%%%
\usepackage{amsthm}

\theoremstyle{definition}
\newtheorem{definition inner}{Definition}[section]
\newenvironment{definition}[1][]
{\begin{definition box}\begin{definition inner}[#1]\hfill\par}
{\end{definition inner}\end{definition box}}

\theoremstyle{plain}
\newtheorem{theorem inner}{Theorem}[section]
\newenvironment{theorem}[1][]
{\begin{theorem box}\begin{theorem inner}[#1]\hfill\par}
{\end{theorem inner}\end{theorem box}}

\theoremstyle{definition}
\newtheorem*{proof inner}{\textit{Proof}}
\renewenvironment{proof}[1][]
{\begin{proof inner}[#1]\hfill\par}
{\qed\end{proof inner}}

\theoremstyle{remark}
\newtheorem*{remark}{\textit{Remark}}
%%%%%%%%%%%%%%%%%%%%%%%%%%%%%%%%%%%%%%%%%%%%%%


\newcommand{\mycases}[2][0.8]{
    \left\{\vphantom{\scalebox{#1}{
        $\begin{matrix}#2\end{matrix}$
    }}\right.\!\!\!\!
    \begin{array}{l@{\quad}l}#2\end{array}
}

\newcommand{\ju}{\mathrm{j}}  % 虚数单位
\newcommand{\e}{\mathrm{e}{\mkern 1 mu}}

\newcommand{\Ev}{\,\mathcal{E}\ensuremath{v}\qty}
\newcommand{\Od}{\,\mathcal{O}\ensuremath{d}\qty}

\renewcommand{\Re}{\,\mathcal{R}\ensuremath{e}\qty}
\renewcommand{\Im}{\,\mathcal{I}\ensuremath{m}\qty}


\begin{document}

\part*{Sampling}

\section{Representation by Samples}

\begin{theorem}[Nyquist-Shannon Sampling Theorem]
    Let $x(t)$ be a \textbf{band-limited} signal
    with a maximum frequency component of $\omega_M$, i.e.
    $X(\ju\omega) = 0$ for all $\abs{\omega} > \omega_M$.
    Then, $x(t)$ can be completely determined by its samples $x[n] = x(n T)$,
    where $T$ is the sampling period,
    if the sampling frequency $\omega_s=2\pi/T$ satisfies the condition:
    \begin{equation*}
        \omega_s > 2\omega_M
    \end{equation*}
    where $\omega_M$ is then called the \textbf{Nyquist frequency},
    and $2\omega_M$ is the \textbf{Nyquist rate}.
\end{theorem}

\subsection{Impulse Train Sampling}

\textbf{Sampling Function}: \quad
$\displaystyle p(t) = \sum_{n=-\infty}^{\infty} \delta(t-n T_s)$
\qquad with \textbf{sampling period} $T_s$
and \textbf{sampling frequency} $\omega_s = \dfrac{2\pi}{T_s}$.
\begin{align*}
    x_p(t) &= x(t) p(t) \\
    &= \sum_{n=-\infty}^{\infty} x(n T_s) \delta(t-n T_s)
\end{align*}
\begin{align*}
    p(t) = \sum_{n=-\infty}^{\infty} \delta(t-n T_s)
    \quad\xrightleftharpoons[\mathcal{F}^{-1}]{\mathcal{F}}\quad
    P(\ju\omega) &= \frac{2\pi}{T_s} \sum_{k=-\infty}^{\infty} \delta\qty(\omega - k \omega_s)
\end{align*}
\begin{align*}
    x_p(t) = x(t) p(t)
    \quad\xrightleftharpoons[\mathcal{F}^{-1}]{\mathcal{F}}\quad
    X_p(\ju\omega) &= \frac{1}{2\pi} X(\ju\omega) \ast P(\ju\omega) \\
    &= \frac{1}{T_s} \sum_{k=-\infty}^{\infty} X\qty(\ju\qty(\omega - k \omega_s)) \\
\end{align*}
For a periodic signal $x(t)$ whose Fourier series is $a_k$,
\begin{gather*}
    x(t)
    \quad\xrightleftharpoons[\mathcal{FS}^{-1}]{\mathcal{FS}}\quad
    a_k \\
    p(t) = \sum_{n=-\infty}^{\infty} \delta(t-n T_s)
    \quad\xrightleftharpoons[\mathcal{FS}^{-1}]{\mathcal{FS}}\quad
    b_k = \frac{1}{T_s} \\
    x_p(t) = x(t) p(t)
    \quad\xrightleftharpoons[\mathcal{FS}^{-1}]{\mathcal{FS}}\quad
    c_k = \frac{1}{T_s} \sum_{n=-\infty}^{\infty} a_{k - n \cdot \frac{T}{T_s}}
\end{gather*}

\subsection{}
Suppose $x(t)$ is a \textbf{band-limited} signal with
a maximum frequency component of $\omega_M$.
\begin{align*}
    x(t)
    \quad&\xrightleftharpoons[\mathcal{F}^{-1}]{\mathcal{F}}\quad
    X(\ju\omega) \\
    x_p(t) := \sum_{n=-\infty}^{\infty} x(n T_s) \delta(t-n T_s)
    \quad&\xrightleftharpoons[\mathcal{F}^{-1}]{\mathcal{F}}\quad
    X_p(\ju\omega) =
    \frac{1}{T_s} \sum_{k=-\infty}^{\infty} X\qty(\ju\qty(\omega - k \omega_s))
    \qq{where} \omega_s = \frac{2\pi}{T_s} \\
    x_d[n] := x(n T_s)
    \quad&\xrightleftharpoons[\mathcal{F}^{-1}]{\mathcal{F}}\quad
    X_d(\e^{\ju\omega}) = X_p\qty(\ju\frac{\omega}{T_s})
    = \frac{1}{T_s} \sum_{k=-\infty}^{\infty} X\qty(\ju\frac{\omega-2 k\pi}{T_s})
\end{align*}


\section{Reconstruction via Interpolation}
Impulse train sampling is not practical.

\subsection{Zero-Order Hold (ZOH)}

\subsection{First-Order Hold (FOH)}

\section{Aliasing}


\section{DT processing of CT signals}


\section{Sampling of DT signals}



\end{document}